\chapter{第十八章：系统设计面试}

AI Agent 岗位的系统设计面试，和传统后端系统设计最大的不同，在于你不仅要设计`\texttt{服务''，还要设计}\texttt{模型行为控制层''。面试官会特别关注：检索链路是否可靠、工具调用是否安全、上下文是否可管理、成本是否能压住、评估是否能闭环。本章给出 6 个高频案例，每个案例都按}`需求---架构---数据流---扩展---成本---权衡''的答题模板展开。


在真正的面试里，系统设计题往往不是要求你把每个组件都讲成论文，而是看你能否抓住`\texttt{主链路''和}\texttt{风险点''。一个成熟的回答通常具备四个特征：第一，先澄清成功指标，比如准确率优先还是时延优先；第二，主动拆出在线链路和离线链路，因为 AI 系统很多能力提升都发生在离线评估与数据回流；第三，明确模型只是系统中的一个组件，真正决定生产可用性的往往是权限、缓存、回退、监控和人工兜底；第四，能说出为什么不用某些更}\texttt{酷''的方案。你可以把本章所有案例都理解成同一个练习：如何把}`模型能力''包进一个可靠的软件系统。


\bigskip\hrule\bigskip

\section{案例一：企业知识库 QA 系统}

\subsection{1. 需求分析}
\textbf{功能性需求}：支持 PDF/Word/HTML/Confluence 摄取；支持多轮问答；支持文档级权限控制；回答必须给出引用来源；支持用户反馈``有用/无用''。


\textbf{非功能性需求}：P95 延迟控制在 3 到 6 秒；权限不能串数据；检索结果可解释；文档更新后 10 分钟内可检索；支持千人并发与百万级 chunk。


\subsection{2. 高层架构}
\begin{codeblock}
        [Doc Sources]
 PDF  Word  HTML  Confluence
          |
          v
   [Ingestion Pipeline]
 OCR / Parse / Clean / Chunk
          |
          +--> [Metadata DB]
          |
          +--> [Sparse Index(BM25)]
          |
          +--> [Vector Index]

User -> API Gateway -> Auth -> Query Rewrite
                          |
                          v
               Hybrid Retrieval -> Reranker
                          |
                          v
                 Context Builder + ACL Filter
                          |
                          v
                  LLM Answer Generator
                          |
                 Citation + Feedback Logger
\end{codeblock}

\subsection{3. 组件深入}
\begin{itemize}[leftmargin=2em]
\item \textbf{摄取层}：按文档类型分别走解析器。PDF 需要 OCR 和版面恢复；Confluence 则保留页面层级、作者、更新时间。
\item \textbf{切块层}：优先按标题、段落、列表、代码块做结构化切分，再做长度约束，避免纯固定字符切块。
\item \textbf{索引层}：稀疏索引处理关键词精确匹配，向量索引处理语义召回，metadata 里保存文档 ID、部门、密级、租户、更新时间。
\item \textbf{权限层}：ACL 既要在召回前过滤，也要在召回后复核，防止 embedding 召回到越权文档。
\item \textbf{生成层}：提示词强制要求``仅基于证据作答，缺证据则明确说明''。
\item \textbf{反馈层}：记录 query、召回文档、最终答案、用户反馈，供离线评估和 hard negative mining 使用。
\end{itemize}

\subsection{4. 数据流}
\begin{enumerate}[leftmargin=2em]
\item 文档进入 ingestion pipeline。
\item 解析文本、图表、附件，生成 chunk 和 metadata。
\item 同步写入 BM25 索引、向量索引与元数据库。
\item 用户提问后，经认证拿到用户身份和权限标签。
\item query rewrite 生成检索用问题；hybrid retrieval 并行召回。
\item ACL filter 去掉无权限 chunk，reranker 做精排。
\item context builder 压缩上下文，附带来源页码。
\item LLM 生成答案并输出 citation。
\item 用户反馈回流到评估系统。
\end{enumerate}

\subsection{5. 扩展性考虑}
\begin{itemize}[leftmargin=2em]
\item 解析服务异步化，防止大文件阻塞。
\item 索引按租户和业务线分片，减少过滤开销。
\item 热门 query 可缓存最终答案和检索结果。
\item 文档更新走增量重建，避免全量索引重算。
\item 对长文档引入 parent-child retrieval：召回小块，回溯到大段上下文。
\end{itemize}

\subsection{6. 成本估算}
以 100 万 chunk、每日 5 万查询为例：

\begin{itemize}[leftmargin=2em]
\item 向量存储和 BM25 索引是固定成本。
\item LLM 调用成本由 top-k、chunk 长度、是否使用 reranker 决定。
\item 若每次召回 8 个 chunk、平均 400 token，每日仅上下文输入就可能达到上亿 token。
\end{itemize}
优化顺序通常是：先优化 chunk 和重排，再考虑更便宜模型或缓存。


\subsection{7. 关键权衡}
\begin{itemize}[leftmargin=2em]
\item \textbf{混合检索 vs 纯向量检索}：前者更稳，后者更简单。
\item \textbf{强权限过滤 vs 高召回}：权限越严格，召回越可能下降，需要索引分区设计配合。
\item \textbf{长上下文直接塞模型 vs 检索压缩}：前者开发快，后者长期成本更优。
\end{itemize}

\subsection{8. 面试追问与答题抓手}
这类题面试官很喜欢继续追问三件事。第一，\textbf{文档更新一致性}：你要能回答`\texttt{更新后多久生效、旧索引如何回收、失败任务如何补偿''。第二，**权限安全**：你要明确 ACL 不能只在前端做，也不能只在最终展示时做，而应贯穿召回、重排和引用渲染。第三，**效果评估**：不能只说}\texttt{用户觉得好''，而要说召回率、引用准确率、拒答率和用户采纳率。答题时如果你能主动说}`这个系统上线后我会先做灰度，观察高频 query 的引用正确率''，会非常像真正带过项目的人。


\bigskip\hrule\bigskip

\section{案例二：多轮对话智能客服 Agent}

\subsection{1. 需求分析}
\textbf{功能性需求}：识别用户意图；抽取槽位（slot）；调用订单查询、退款、发券、升级人工等工具；保持多轮上下文；支持人工接管；记录服务质量。


\textbf{非功能性需求}：高峰期秒级响应；高风险操作可审计；退款操作必须幂等；多语言支持；故障时能安全降级到 FAQ 或人工。


\subsection{2. 高层架构}
\begin{codeblock}
User -> Channel(Web/App/WhatsApp)
          |
          v
      Session Router
          |
   +------+------+
   | Intent Model |
   | Policy Layer |
   +------+------+
          |
          v
      Agent Orchestrator
   /       |         \
Order API Refund API Escalation API
   \       |         /
          v
     Response Composer
          |
          v
 Human Handoff / QA Monitor / Memory Store
\end{codeblock}

\subsection{3. 组件深入}
\begin{itemize}[leftmargin=2em]
\item \textbf{意图识别}：可用轻量分类模型或小型 LLM，把`\texttt{查物流''}\texttt{催发货''}`申请退款''先分流，减少主模型成本。
\item \textbf{槽位填充}：订单号、退款原因、商品 ID、联系电话等信息需要结构化抽取，并支持追问补齐。
\item \textbf{Agent 编排器}：决定是继续问用户、直接调用工具、还是转人工。
\item \textbf{上下文管理}：短期保留最近若干轮，长期只保存用户画像和工单摘要。
\item \textbf{人工接管}：要能把会话摘要、已执行工具、当前情绪、风险标签一并交给人工坐席。
\end{itemize}

\subsection{4. 数据流}
\begin{enumerate}[leftmargin=2em]
\item 用户消息进入会话路由层。
\item 系统读取 session state 和 CRM 基础信息。
\item 意图识别器输出 intent 和置信度。
\item 槽位提取器判断参数是否齐全，不齐则追问。
\item 编排器根据 policy 决定调用工具或升级人工。
\item 工具结果写回 conversation state。
\item LLM 生成自然语言回复。
\item 会话结束后写入摘要、满意度和质检标签。
\end{enumerate}

\subsection{5. 扩展性考虑}
\begin{itemize}[leftmargin=2em]
\item 使用异步工具执行，避免慢接口拖垮整体体验。
\item 将`\texttt{意图识别''和}`开放问答''拆成两层，减少大模型调用频次。
\item 对高频动作如``查订单''做结果缓存。
\item 将监控拆成实时告警和离线质检两条链路。
\end{itemize}

\subsection{6. 成本估算}
客服系统的核心成本不是单次对话，而是日活乘以多轮消息数。节流策略包括：

\begin{itemize}[leftmargin=2em]
\item 首轮先走轻量路由模型。
\item 明确命中 FAQ 时不用主 LLM。
\item 工具结果模板化输出，减少自由生成 token。
\item 长会话做摘要压缩，防止上下文无限增长。
\end{itemize}

\subsection{7. 关键权衡}
\begin{itemize}[leftmargin=2em]
\item \textbf{强规则流程 vs 高自由对话}：规则强则稳定，自由高则体验好。
\item \textbf{全靠 LLM 做意图识别 vs 分层模型}：全靠 LLM 简单，但成本高且不稳定。
\item \textbf{自动退款 vs 人工审批}：越自动越省人力，但风险和审计要求更高。
\end{itemize}

\subsection{8. 面试追问与答题抓手}
客服题的高分关键，是你能否把`\texttt{对话体验''和}\texttt{交易安全''同时讲清楚。很多候选人只会讲 agent loop，却不讲退款幂等、人工接管、情绪识别误判成本，这会显得非常产品化但不工程化。更好的说法是：低风险问题尽量自动化，高风险动作必须提升置信门槛，并保留人工兜底。若被问到}`如何减少 token 成本''，你可以回答：先用小模型做路由，再用模板化工具回复覆盖大部分事务型请求，只有复杂解释类问题才进入主模型。


\bigskip\hrule\bigskip

\section{案例三：AI 代码审查助手}

\subsection{1. 需求分析}
\textbf{功能性需求}：接收 Git PR webhook；拉取 diff 和上下文代码；生成 review comment；支持按严重程度分类；支持开发者反馈``误报/采纳''。


\textbf{非功能性需求}：误报率要低；不能泄露私有代码；对大 PR 也能在可接受时间内完成；需要可审计；与 GitHub/GitLab 易集成。


\subsection{2. 高层架构}
\begin{codeblock}
Git PR Webhook
      |
      v
  Review Ingestor -> Diff Parser -> File Classifier
                              |          |
                              |          +--> Rule Engine
                              v
                        Context Retriever
                              |
                              v
                      LLM Review Generator
                              |
                      Severity Calibrator
                              |
                  Comment Publisher / Feedback Store
\end{codeblock}

\subsection{3. 组件深入}
\begin{itemize}[leftmargin=2em]
\item \textbf{Diff Parser}：不仅要看新增/删除行，还要定位函数、类、依赖调用链。
\item \textbf{Rule Engine}：先用静态规则拦截明显问题，如密钥提交、空指针风险、缺少测试，再把复杂语义问题交给 LLM。
\item \textbf{Context Retriever}：为每个 diff hunk 补充相邻函数、历史实现、相关测试文件、相似历史缺陷。
\item \textbf{Review Generator}：提示词应要求``只报高置信、可行动的问题''，否则误报会严重伤害体验。
\item \textbf{Feedback Store}：记录开发者是否采纳评论，用于后续 prompt 和排序优化。
\end{itemize}

\subsection{4. 数据流}
\begin{enumerate}[leftmargin=2em]
\item PR 事件触发系统。
\item 拉取 diff、文件树、历史提交和相关测试。
\item 文件分类器判断是后端、前端、基础设施还是配置文件。
\item 规则引擎做第一轮过滤。
\item 上下文检索器补足函数级或模块级上下文。
\item LLM 生成候选 review comments。
\item 严重级别校准器过滤低置信评论。
\item 发布到 PR，并记录后续开发者反馈。
\end{enumerate}

\subsection{5. 扩展性考虑}
\begin{itemize}[leftmargin=2em]
\item 大 PR 可按文件并行处理，再在最后聚合。
\item 高频仓库可做 embedding 索引，缓存常见模块语义。
\item 对低风险文件（如锁文件、生成文件）直接跳过。
\item 评论去重很重要，避免多个 agent 重复指出同一问题。
\end{itemize}

\subsection{6. 成本估算}
成本主要来自三部分：拉代码上下文、模型推理、反馈存储。对大仓库来说，真正贵的是``把足够上下文喂给模型''。因此实践中常用分级策略：先规则筛一遍，只让少量高风险片段进入大模型；对轻量文件使用便宜模型，对核心业务逻辑使用强模型。


\subsection{7. 关键权衡}
\begin{itemize}[leftmargin=2em]
\item \textbf{高召回 vs 低误报}：代码审查场景更怕误报，宁可少报，也别天天打扰开发者。
\item \textbf{文件级独立分析 vs 全局依赖理解}：前者快，后者更准但更贵。
\item \textbf{在线实时评论 vs 批量总结}：实时感知更强，批量更省 token。
\end{itemize}

\subsection{8. 面试追问与答题抓手}
代码审查题常见追问是：`\texttt{你怎么证明这东西不会天天乱报？''此时不要只说}\texttt{换更强模型''，而要强调分层控制：规则优先、上下文补全、置信度过滤、开发者反馈回流。你还可以补一句非常加分的话：真正好的 AI review assistant 不应追求覆盖所有问题，而应专注于高价值、高置信、可行动的问题，比如安全漏洞、并发缺陷、事务边界错误、权限判断遗漏。这个回答会明显比}`帮开发者提代码风格建议''更成熟。


\bigskip\hrule\bigskip

\section{案例四：金融研报分析 Agent}

\subsection{1. 需求分析}
\textbf{功能性需求}：摄取 SEC 文件、新闻、财务 API、内部研报；抽取指标；生成趋势分析；采用 researcher、analyst、writer 多 agent 协作；输出带证据和置信度的报告。


\textbf{非功能性需求}：高准确率、可追溯、时效性强、合规审计严格；输出不能把推测写成事实；需要版本控制和复盘能力。


\subsection{2. 高层架构}
\begin{codeblock}
SEC/News/API/Internal Reports
            |
            v
      Data Normalization Layer
            |
   +--------+--------+
   |                 |
Structured Store   Document Store
   |                 |
   +--------+--------+
            v
      Researcher Agent
            |
      Analyst Agent
            |
       Writer Agent
            |
  Compliance Checker + Citation Verifier
            |
            v
         Final Report
\end{codeblock}

\subsection{3. 组件深入}
\begin{itemize}[leftmargin=2em]
\item \textbf{数据标准化层}：统一公司名称、ticker、时间粒度、币种和来源标签。
\item \textbf{Researcher Agent}：负责检索最新事实，整理证据包。
\item \textbf{Analyst Agent}：基于结构化指标和证据做趋势推断、横向比较和风险识别。
\item \textbf{Writer Agent}：把结论整理成可读报告，但不得发明新事实。
\item \textbf{Compliance Checker}：检查是否缺引用、是否出现夸张表述、是否混淆历史事实与预测。
\end{itemize}

\subsection{4. 数据流}
\begin{enumerate}[leftmargin=2em]
\item 多源数据进入标准化层。
\item 文档数据进入检索索引，结构化数据进入分析仓库。
\item Researcher agent 先收集证据，形成 research brief。
\item Analyst agent 基于证据和数值数据生成分析框架。
\item Writer agent 输出研报草稿。
\item 合规检查器校验事实、引用和风险措辞。
\item 通过后入库，并记录版本与模型配置。
\end{enumerate}

\subsection{5. 扩展性考虑}
\begin{itemize}[leftmargin=2em]
\item 新闻与公告需要流式摄取，支持分钟级更新。
\item 结构化指标可预聚合，降低每次问答的计算量。
\item 不同市场和行业可分域部署检索索引。
\item 高价值报告应支持人工编辑和二次审核。
\end{itemize}

\subsection{6. 成本估算}
金融场景的 LLM 成本通常不是最高，最高的是`\texttt{高质量数据源''和}`人工审核''。模型成本可以通过三段式结构控制：research 用较便宜模型初筛，analysis 与 final writing 用更强模型；对重复请求复用 research brief 和缓存统计指标。


\subsection{7. 关键权衡}
\begin{itemize}[leftmargin=2em]
\item \textbf{多 agent 专业分工 vs 单 agent 端到端}：前者更可控，后者更简单。
\item \textbf{自动总结 vs 人工复核}：金融领域高风险，最后一道人工审批通常不可省。
\item \textbf{时效性 vs 准确性}：越追求实时，越要防止引用未验证新闻。
\end{itemize}

\subsection{8. 面试追问与答题抓手}
金融题最容易被追问的是`\texttt{如果新闻和财报冲突怎么办''。推荐的答法是：先保留冲突，而不是强行统一结论；把不同来源、时间、可信度并列展示，再由 analyst agent 给出条件化判断。另一类追问是}\texttt{如何避免模型把推测写成事实''，这时你要主动提出模板化输出：把事实、分析、预测、风险拆成四栏，并要求每一条结论都附来源编号。只要你把}\texttt{合规''和}`可追溯''反复强调，金融场景的答题质感会立刻提升。


\bigskip\hrule\bigskip

\section{案例五：多 Agent 协作平台}

\subsection{1. 需求分析}
\textbf{功能性需求}：提供 agent 注册与发现；支持任务路由、负载均衡、共享记忆、执行监控、成本归因；支持不同团队接入自己的 agent。


\textbf{非功能性需求}：平台级多租户隔离；高可用；审计完备；延迟可接受；支持异步任务和长任务。


\subsection{2. 高层架构}
\begin{codeblock}
                +-------------------+
Client/API ---> | Gateway & Auth    |
                +---------+---------+
                          |
                    Task Router
          +---------------+----------------+
          |                                |
   Agent Registry                    Workflow Engine
          |                                |
  Capability Index                   State / Queue / Retry
          |                                |
   +------+------+                   +-----+------+
   |             |                   |            |
Planner Agent  Tool Agent        Memory Store  Event Bus
   |             |                   |            |
   +-------------+-------------------+------------+
                          |
                    Monitor & Billing
\end{codeblock}

\subsection{3. 组件深入}
\begin{itemize}[leftmargin=2em]
\item \textbf{Agent Registry}：登记 agent 的能力、输入 schema、工具权限、SLA、所属租户。
\item \textbf{Task Router}：根据任务类型、预算、优先级和历史表现选择 agent。
\item \textbf{Workflow Engine}：管理任务状态、超时、重试、补偿、人工审批。
\item \textbf{Shared Memory}：保存任务上下文、共享文档、阶段摘要，但要按租户隔离。
\item \textbf{Billing / Cost Allocation}：记录每次 token、工具调用、GPU 时间、外部 API 费用，支持按团队回收成本。
\end{itemize}

\subsection{4. 数据流}
\begin{enumerate}[leftmargin=2em]
\item 客户端提交任务与预算约束。
\item 网关完成鉴权并写入 trace id。
\item 路由器查询 registry 选择最合适 agent。
\item workflow engine 启动任务，必要时拆成多个子任务。
\item agent 执行过程中从 shared memory 读取上下文并回写结果。
\item event bus 推送监控事件到 dashboard。
\item billing 模块聚合成本，形成账单和优化报表。
\end{enumerate}

\subsection{5. 扩展性考虑}
\begin{itemize}[leftmargin=2em]
\item Registry 需要 capability index，避免线性扫描全部 agent。
\item Workflow engine 应支持幂等恢复，避免长任务中途失败全盘重来。
\item Shared memory 可冷热分层：热数据放 KV/Redis，冷数据放对象存储。
\item 任务队列按优先级和租户权重调度，防止单租户打爆平台。
\end{itemize}

\subsection{6. 成本估算}
平台类系统的成本来自三层：控制面服务、模型调用、外部工具。最容易失控的是``无边界的 agent 递归调用''。因此要从平台层限制最大深度、最大步数、最大 token 和工具预算，并把成本标签透传到每个子任务。


\subsection{7. 关键权衡}
\begin{itemize}[leftmargin=2em]
\item \textbf{集中式编排 vs 去中心化协作}：集中式更可控，去中心化更灵活。
\item \textbf{共享记忆越多越好？} 不一定，过多共享会带来串扰、泄露和调试困难。
\item \textbf{统一强模型 vs 按任务混合模型}：统一简单，混合更省钱也更灵活。
\end{itemize}

\subsection{8. 面试追问与答题抓手}
平台题的难点是很多候选人会把它答成`\texttt{多个 agent 排队调用''，但真正的平台化设计一定要包含注册、发现、路由、计费、监控、隔离和审计。面试时如果被问}\texttt{如何防止某个 agent 失控''，你可以从三层回答：任务层限制最大递归深度，运行层限制 token 和工具预算，平台层限制租户配额和告警阈值。再补一句}`平台的价值不是让 agent 更多，而是让 agent 可治理''，基本就是很强的答案了。


\bigskip\hrule\bigskip

\section{案例六：实时 RAG 搜索引擎}

\subsection{1. 需求分析}
\textbf{功能性需求}：持续抓取网页；实时更新索引；支持 query understanding、query reformulation、混合检索、流式生成答案、来源引用与事实验证。


\textbf{非功能性需求}：分钟级索引延迟；高 QPS；高可用；结果新鲜度高；要防网页注入和恶意 SEO 污染。


\subsection{2. 高层架构}
\begin{codeblock}
Web Crawl -> Parse/Clean -> Dedup -> Index Builder
    |             |            |          |
    |             |            |          +--> BM25 Index
    |             |            |          +--> Vector Index
    |             |            |
    |             |            +--> Trust / Spam Score
    v
 Scheduler

User Query -> Query Understanding -> Reformulation
                              |
                              v
                  Hybrid Retrieval + Freshness Boost
                              |
                              v
                    Verification / Reranking
                              |
                              v
                   Streaming Answer Generator
                              |
                              v
                      Citation Renderer
\end{codeblock}

\subsection{3. 组件深入}
\begin{itemize}[leftmargin=2em]
\item \textbf{抓取层}：支持增量 crawl、robots 策略、站点优先级和去重。
\item \textbf{解析层}：提取主正文、标题、时间、作者、链接关系，过滤广告和导航垃圾。
\item \textbf{质量层}：给页面打可信度分数、垃圾站点分数、时效分数。
\item \textbf{查询理解层}：识别用户是找定义、最新新闻、比较分析还是导航型查询。
\item \textbf{验证层}：对高风险结论做多源交叉验证，尽量避免单页单源定结论。
\end{itemize}

\subsection{4. 数据流}
\begin{enumerate}[leftmargin=2em]
\item 爬虫抓取网页，解析器抽正文和 metadata。
\item 去重与垃圾检测后进入索引构建。
\item 用户查询进入 query understanding，必要时改写出多个子查询。
\item hybrid retrieval 同时按语义和关键词召回，并叠加新鲜度权重。
\item reranker 结合相关性、可信度和时效性排序。
\item answer generator 边生成边流式返回，并插入 citation 标记。
\item 最终界面支持点击原文和反馈纠错。
\end{enumerate}

\subsection{5. 扩展性考虑}
\begin{itemize}[leftmargin=2em]
\item 抓取和索引构建需要完全异步化。
\item 对热门站点做差分更新，而不是全站重抓。
\item 检索层可用分片索引与多级缓存。
\item 流式生成要支持客户端取消，避免无效 token 成本。
\item 高风险主题加入更强验证或直接只返回链接摘要。
\end{itemize}

\subsection{6. 成本估算}
实时搜索的成本结构通常是：抓取带宽与存储 > 索引构建 > 检索 > 生成。若追求极高时效，索引更新频率会显著抬升成本。常见优化是：热门 query 结果缓存、对低价值页面降采样、对简单导航型查询直接返回链接而不触发生成。


\subsection{7. 关键权衡}
\begin{itemize}[leftmargin=2em]
\item \textbf{新鲜度 vs 稳定性}：更新越快，污染数据进入系统的概率也越高。
\item \textbf{直接答案 vs 搜索结果页}：直接答案体验好，但事实风险更大。
\item \textbf{全量生成 vs 结果分层}：简单 query 不必每次都走 LLM。
\end{itemize}

\subsection{8. 面试追问与答题抓手}
实时搜索题经常会被追问`\texttt{为什么不直接像传统搜索那样返回链接''。你的回答可以是：当查询是定义类、比较类、总结类时，生成式答案能显著降低用户理解成本；但对于导航类和高风险查询，应该退回结果页或摘要页。另一种高频追问是}`如何对抗垃圾网页和 SEO 污染''，这时要主动提到可信度评分、多源交叉验证、域名白名单和新站冷启动策略。只讲检索、不讲信息污染治理，在这类题里会丢很多分。


\bigskip\hrule\bigskip

\section{本章答题模板总结}

在系统设计面试中，你可以固定按以下顺序展开：


\begin{enumerate}[leftmargin=2em]
\item \textbf{澄清需求}：用户是谁、核心成功指标是什么。
\item \textbf{拆功能模块}：摄取、检索、生成、工具、记忆、评估、安全。
\item \textbf{画高层图}：哪怕是 ASCII 图，也能显著提升表达清晰度。
\item \textbf{讲数据流}：从输入到输出，逐步解释每个环节。
\item \textbf{补可观测性与安全}：日志、评估、权限、审计、回滚。
\item \textbf{做成本与权衡}：让面试官看到你不是只会``堆组件''。
\end{enumerate}

如果你把本章 6 个案例都练熟，系统设计面试里绝大多数 AI Agent 题目都能被你映射到这些模板上。下一章我们进入编程面试与实操，重点解决`\texttt{现场写代码''和}`take-home 项目''两大痛点。


最后再强调一个实战技巧：系统设计面试不是比谁说得快，而是比谁更会`\texttt{画边界''。当你讲清楚哪些模块必须在线、哪些可以离线，哪些地方必须强一致、哪些地方可以最终一致，哪些风险必须人工兜底、哪些可以自动恢复时，面试官自然会把你看成具备架构判断力的候选人。这种判断力，正是 AI Agent 岗位比普通}`会调模型 API''更值钱的地方。


\section{面试现场表达模板：把系统讲得像你真的做过}

很多人明明会设计，却在表达上吃亏。这里给你一个非常实用的口述模板，几乎可以套到所有 AI Agent 系统设计题里。第一步，先用一句话定义目标，例如``这是一个面向企业内部员工的知识问答系统，核心指标是回答准确率、权限隔离和 P95 延迟''。第二步，立刻拆在线与离线：离线负责数据摄取、清洗、索引、评估；在线负责路由、检索、生成、缓存、审计。第三步，说明最关键的状态对象是什么，比如会话状态、索引版本、权限标签、工具预算。第四步，主动指出最脆弱的环节：RAG 题通常是检索质量和引用正确性，客服题通常是工具安全和人工接管，平台题通常是编排复杂度和成本失控。第五步，补监控和回滚：哪些指标异常会触发降级、哪些版本切换支持灰度、哪些场景必须人工审批。


如果面试官中途打断，你也不要慌。系统设计面试很少要求你从头到尾一口气说完，更常见的是对某个点深入追问。你的目标不是背完整稿子，而是让每个模块都能单独展开 1 到 2 分钟。换句话说，本章六个案例最有价值的地方，不只是给你六个答案，而是帮你练出一种``被追问也不散架''的表达结构。


还有三个常见禁忌值得你记住。第一，不要把系统设计题答成`\texttt{模型选型题''，只谈 GPT、Claude、Gemini 谁强谁弱，却不谈缓存、回退、审计、权限和评估。第二，不要把所有问题都推给}`大模型更强版本''，因为面试官想看到的是架构判断，不是预算无限的幻想。第三，不要忽略失败路径：真实系统不是只在 happy path 上运行，索引会延迟、工具会超时、文档会解析失败、用户会越权访问，这些都应该在你的回答里有明确归宿。只要你能避开这三类常见失误，系统设计面的稳定性就会提升一个层级。


如果你想在收尾时再加一句有力量的话，可以使用这样的模板：``我设计这个系统时，优先保证的是正确性和可治理性，然后再通过缓存、分层模型和异步化逐步优化成本与时延。'' 这句话之所以有效，是因为它把 AI Agent 系统最核心的三组矛盾------准确率、成本、速度------同时点出来了。很多系统设计题没有唯一正确答案，但几乎所有高分答案，都会体现出你知道先保什么、后优化什么。


最后，别忘了系统设计面试也是沟通面试。你画图时越清楚、分层越明确、命名越稳定，面试官越容易跟上你的节奏。很多人方案并不差，只是讲得像临时拼起来的白板笔记；而高分候选人通常会反复使用统一术语，例如``摄取层、索引层、编排层、评估层、治理层''，让整个答案听起来像一个真正维护过的系统。
